\section{Transactional Control Algebra}
\label{sec:algebra}

\subsection{Reachable updates}

Let $S$ be the current state, $\tau$ a transaction, and
$\Delta S^\star(S,\tau)$ its target update.  A controller class
$\mathcal{C}$ induces a reachable set
$\mathcal{R}_{\mathcal{C}}(S,\tau)$.  State-space regret measures distance to
that set, whereas behavioral regret evaluates the best reachable update after
the task readout $\rho$:
\begin{align}
R_{\mathrm{state}}
  &= \min_{\Delta S\in\mathcal{R}_{\mathcal{C}}}
     \lVert\Delta S-\Delta S^\star\rVert^2,\\
R_{\mathrm{beh}}
  &= \min_{\Delta S\in\mathcal{R}_{\mathcal{C}}}
     \lVert\rho(S+\Delta S)-\rho(S+\Delta S^\star)\rVert^2.
\end{align}
The distinction matters whenever state components differ in behavioral
relevance.

\subsection{Control rank}

For descriptor-conditioned target operators $A(z)$, a rank-$r$ controller can
only reach a rank-constrained family.  We separate intrinsic target rank,
learned controller rank, the per-example best-rank floor, and learned
excess-over-floor.  The prospective E10b rule asks for the smallest learned
rank that satisfies a floor-aware quality condition; it does not equate an
oracle upper bound with learned recovery.

\subsection{Shared bases and demand algebra}

A shared diagonal controller represents a demand family as
\begin{equation}
  A_\tau = Q\,\mathrm{Diag}(c_\tau)\,Q^\top
\end{equation}
for one shared basis $Q$.  Axis-aligned and common-rotated commuting families
admit this form.  Transaction-dependent noncommuting families need not.
Representation learning can absorb a common rotation, but one shared basis
does not thereby represent arbitrary noncommutativity.  This connects the
controller class to simultaneous diagonalization
\citep{bunsegerstner1993simultaneous,cardoso1996jacobi,glashoff2013almost}.

\subsection{Nested control freedoms}

The controlled lattice is
\begin{equation}
\text{tied scalar}
\subset \text{dual scalar}
\subset \text{diagonal value}
\subset \text{separate address}
\subset \text{state-aware}.
\end{equation}
Each inclusion adds a named freedom.  Our tests ask whether an adjacent
inclusion helps on the matched demand while satisfying registered
simpler-demand and unaffected-retention guardrails.
